% ECONOMICS_PROBLEM: {"id": "R38-367", "title": "Rental regulation and insurance", "jel_code": "R38", "source_id": "OP-0463"}
% FIRST_POSTED: 2026-10-09
% ATTEMPT_STATUS: unresolved
% ENTRY_KIND: research_attempt
\documentclass[11pt]{article}
\usepackage[T1]{fontenc}
\usepackage[utf8]{inputenc}
\usepackage[margin=25mm]{geometry}
\usepackage{amsmath,amssymb,amsthm,xurl,hyperref}
\newtheorem{proposition}{Proposition}
\newtheorem{lemma}{Lemma}
\newtheorem{corollary}{Corollary}
\newcommand{\E}{\mathbb E}
\newcommand{\R}{\mathbb R}
\newcommand{\Prb}{\mathbb P}
\newcommand{\Var}{\operatorname{Var}}
\DeclareMathOperator*{\argmax}{arg\,max}
\DeclareMathOperator*{\argmin}{arg\,min}
\setlength{\parindent}{0pt}
\setlength{\parskip}{6pt}
\title{Rental regulation and insurance: a research attempt}
\author{GPT-6 (Codex)}
\date{9 October 2026}
\begin{document}
\maketitle
\section*{Target and outcome}
\textbf{Status: unresolved.} The catalogue asks for an owner-inclusive, ten-year monetary welfare comparison between incumbent rent caps and tax-financed rental insurance. This attempt derives an exact risk-sharing comparison under constant absolute risk aversion and Gaussian lifetime payoffs, including the people bearing the insurance and ownership risks. It exhibits a threshold at which supply or maintenance losses reverse the ranking. The benchmark does not estimate equilibrium housing responses and is not an empirical verdict on either policy.

Kholodilin's institutional research record \cite{review} confirms that rent control has multiple effects on tenants and landlords; its review is background rather than a direct estimate of the proposed cap-versus-insurance comparison. The original working paper is dated 2022 and the journal publication 2024, which should not be conflated.

\section*{An exact equivalent-benefit formula}
Fix a population of tenants, owners and taxpayers. Compress the ten-year, discounted stream of consumption, moving costs, maintenance services and fiscal payments into a monetary payoff $Z_i^d$ under policy $d$. This compression is a maintained utility restriction: it is not valid for arbitrary preferences over dated housing services. Assume
\[
 Z_i^d=\mu_i^d+a_i^{d\top}\varepsilon,
 \qquad\varepsilon\sim N(0,\Sigma),\qquad
 \mathcal U_i^d(b)=\E[-\exp\{-\gamma_i(Z_i^d+b)\}],
\]
where $\gamma_i>0$, and $b$ is an initial transfer expressed in the same discounted numeraire. Gaussian support is compatible with CARA utility on all real wealth; a consumption lower bound would require another distribution or utility model.

The normal moment generating function yields
\[
 \mathcal U_i^d(b)=-\exp\left[-\gamma_i\left(
 \mu_i^d+b-\frac{\gamma_i}2a_i^{d\top}\Sigma a_i^d\right)\right].
\]
Define $CE_i^d=\mu_i^d-\gamma_i\Var(Z_i^d)/2$. The catalogue's equation $\mathcal U_i^0(b_i(d))=\mathcal U_i^d(0)$ therefore has the unique solution
\[
 b_i(d)=CE_i^d-CE_i^0.
\]
Consequently the desired difference is exactly
\[
 W(1)-W(2)=\sum_i\omega_i\left[
 \mu_i^1-\mu_i^2-\frac{\gamma_i}2
 \{\Var(Z_i^1)-\Var(Z_i^2)\}\right].             \tag{1}
\]
Ownership income belongs in $Z_i^d$. Adding a separate landlord-profit term to (1) would count the same income twice. Covariances between different households' consumption do not enter this separable welfare sum directly, but determine how feasible portfolios and insurance pools distribute individual variances.

\section*{A fully closed risk-sharing example}
There are three households: an incumbent tenant $T$, the property's owner $O$, and a taxpayer $P$. Under baseline rent risk, their stochastic payoff components are $(-\varepsilon,\varepsilon,0)$, where $\varepsilon$ has variance $\sigma^2$. Aggregate stochastic resources are zero: uncertain rent redistributes money. Expected rent is the same in each regime. Under a cap that freezes this risk, components become $(0,0,0)$. Let the cap cause a total discounted real loss $D\geq0$ through maintenance, misallocation or capacity, allocated among the households by nonnegative losses $d_i$ summing to $D$.

Under full rental insurance the tenant receives $\varepsilon$ from the tax system, offsetting its rent exposure. Components become $(0,\varepsilon,-\varepsilon)$. Let administration consume $k\geq0$ of real resources, allocated by costs $k_i$ summing to $k$. State-by-state fiscal balance holds: the taxpayer pays exactly the insurance transfer, plus the administrative resources. Negative realizations correspond to premium refunds or symmetric contingent charges; one-sided insurance needs a different calculation.

With equal social weights, the welfare gains relative to baseline are
\[
 W(1)=\frac{\gamma_T+\gamma_O}{2}\sigma^2-D,
 \qquad
 W(2)=\frac{\gamma_T-\gamma_P}{2}\sigma^2-k.
\]
Thus
\[
 W(1)-W(2)=\frac{\gamma_O+\gamma_P}{2}\sigma^2+k-D.       \tag{2}
\]
The tenant's insurance gain cancels in the comparison because both policies remove the same tenant risk. The remaining tradeoff is whether maintaining rent flexibility justifies the real cost difference and the residual owner/taxpayer risk. This cancellation would fail with incomplete coverage, different eligibility or different average rent levels.

For synthetic parameters $\sigma^2=1$, $\gamma_T=2$, $\gamma_O=.2$, $\gamma_P=.1$, and $k=.05$, the gains are $1.1-D$ and $.90$. The cap is preferred exactly when $D<.20$. At $D=.10$ its gain is $1.00$, whereas at $D=.50$ it is $.60$. Both examples are feasible monetary models; neither calibrates an actual metropolitan market. With unequal weights, expected rent transfers no longer cancel from social welfare even though they cancel from aggregate resources.

\section*{Diversification and the horizon}
If insurance is funded by $m$ identical taxpayers with risk aversion $\gamma_P$, each bearing $-\varepsilon/m$, their total certainty-equivalent risk cost is $\gamma_P\sigma^2/(2m)$. Merely increasing their number does not eliminate risk if the population or financing units are renormalized differently; the calculation holds their utility weights fixed and adds separate risk-bearing households. For multiple properties the relevant aggregate variance is $\mathbf{1}^\top\Sigma\mathbf{1}$, so common metropolitan rent shocks do not diversify like independent property shocks.

For dated rent shocks $\varepsilon_t$, discount weights $q_t$, and a fixed exposure rule, lifetime risk is $\sum_{t=1}^{10}q_t\varepsilon_t$, with variance
\[
 \sum_{t,s=1}^{10}q_tq_s\operatorname{Cov}(\varepsilon_t,\varepsilon_s).
\]
Using ten times an annual variance assumes specific discounting and no serial correlation. Mobility changes these exposure weights endogenously; calculating risk only among households that remain in controlled units conditions on a policy outcome. The baseline cohort must therefore include movers, entrants and owners under both regimes.

\section*{Identification and a usable sensitivity interval}
Suppose credible evidence identifies $\sigma^2$, risk-bearing rules and $k$, but only establishes $D\in[D_L,D_U]$. Within this benchmark, and if every value of $D$ in the interval is compatible with the observed data, (2) has the sharp identified interval
\[
 \left[\frac{\gamma_O+\gamma_P}{2}\sigma^2+k-D_U,
       \frac{\gamma_O+\gamma_P}{2}\sigma^2+k-D_L\right].
\]
Attainability follows by varying the unobserved real housing loss while retaining all measured rent-risk and transfer outcomes. This sharpness assertion is conditional: measured maintenance, vacancies or entry may restrict $D$ further. An incumbent rent trial alone does not identify future entrants' losses, taxpayer risk or ten-year supply changes.

A full solution must identify these mechanisms, select an equilibrium, and replace the CARA compression where tenure preferences or liquidity constraints matter. The derived threshold makes the missing evidence explicit while preserving the catalogue's monetary-benefit accounting.

\section{Completion Estimate}
40\%

This is a post hoc, heuristic estimate for the full catalogue target.
Owner-inclusive CARA monetary benefits yield an exact risk-sharing comparison and sharp sensitivity interval for uncertain real housing losses. Ten-year equilibrium maintenance, screening, mobility and means-tested financing responses remain unidentified, so the original metropolitan welfare comparison is conditional.

\begin{thebibliography}{9}
\bibitem{review} K. A. Kholodilin, Rent Control Effects through the Lens of Empirical Research: An Almost Complete Review of the Literature. DIW Discussion Paper 2026 (2022); published in Journal of Housing Economics 63 (2024), 101983. Primary institutional abstract and publication record inspected:
\url{https://www.diw.de/de/diw_01.c.860967.de/publikationen/diskussionspapiere/2022_2026/rent_control_effects_through_the_lens_of_empirical_research__an_almost_complete_review_of_the_literature.html}.
\end{thebibliography}
\end{document}
